\section*{अध्याय \ref{ch:b1:isostasy-gravity} --- समस्थिति और गुरुत्व विसंगतियाँ}

\begin{solution}{exo:b1:isostasy-gravity:1}
(क) $\qty{-150}{\milli\gal} = -150\times\qty{e-5}{m.s^{-2}} = \qty{-1.5e-3}{m.s^{-2}}$,
और $\num{1.5e-3}/9.81 = \num{1.53e-4}$ का $g$। (ख) $\qty{1}{\milli\gal} =
\qty{e-5}{m.s^{-2}}$, $g$ का \num{1.02e-6}; $\qty{1}{\micro\gal} =
\qty{e-8}{m.s^{-2}}$, $g$ का \num{1.02e-9}: जो गुरुत्वमापी एक माइक्रोगैल विभेदित करता है, वह गुरुत्व को लगभग एक अरबवें भाग तक मापता है।
\end{solution}

\begin{solution}{exo:b1:isostasy-gravity:2}
(क) भूमध्य रेखा पर दोनों ज्याएँ शून्य हैं: $\gamma = \qty{9.780327}{m.s^{-2}}$। \ang{45} पर, $\sin^2\varphi = 0.5$ और $\sin^2 2\varphi = 1$:
$\gamma = 9.780327\times(1 + 0.0026512 - 0.0000058) = \qty{9.806200}{m.s^{-2}}$।
ध्रुवों पर, $\sin^2\varphi = 1$ और $\sin 2\varphi = 0$:
$\gamma = 9.780327\times1.0053024 = \qty{9.832186}{m.s^{-2}}$।
(ख) $9.832186 - 9.780327 = \qty{0.051859}{m.s^{-2}}$, अर्थात् \qty{5186}{\milli\gal}। (ग) $3392/5186 = 0.65$: अपकेंद्री प्रभाव अंतर का लगभग दो-तिहाई समझाता है; चपटाई, जो ध्रुवों को केंद्र के निकट रखती है, शेष।
\end{solution}

\begin{solution}{exo:b1:isostasy-gravity:3}
$r = h\rho_c/(\rho_m - \rho_c) = 3\times2800/500 = \qty{16.8}{km}$। \omterm{def:b1:layered-earth:layers}{भूपर्पटी} $h + r = \qty{19.8}{km}$ जितनी मोटी है (अर्थात् $h\rho_m/(\rho_m - \rho_c) =
3\times3300/500$)।
\end{solution}

\begin{solution}{exo:b1:isostasy-gravity:4}
(क) $1/0.3086 = \qty{3.24}{m}$। (ख) $0.01/0.3086 = \qty{0.032}{m}$: ऊँचाइयाँ लगभग \qty{3}{cm} तक ज्ञात होनी चाहिए, इसीलिए एक सर्वेक्षण हर स्टेशन की ऊँचाई उपग्रह-स्थिति-निर्धारण या समतलन से दर्ज करता है।
\end{solution}

\begin{solution}{exo:b1:isostasy-gravity:5}
(क) $\rho = 2700\times100/(100 + h)$: \qty{1}{km} के लिए \qty{2673}{kg.m^{-3}}, \qty{3}{km} के लिए \qty{2621}{kg.m^{-3}}, \qty{5}{km} के लिए \qty{2571}{kg.m^{-3}}।
(ख) $\rho = (2700\times100 - 1030\times4)/96 = \num{265880}/96 =
\qty{2770}{kg.m^{-3}}$: समुद्र के नीचे का स्तंभ सघनतर है। (ग) प्रैट का प्रतिरूप स्थिर गहराई पर घनत्व के पार्श्व परिवर्तन माँगता है। महासागर-तल के नीचे ऐसे परिवर्तन विद्यमान हैं: कटक के निकट युवा, गरम \omterm{def:b1:layered-earth:lithosphere}{स्थलमंडल} दूर के पुराने, ठंडे \omterm{def:b1:layered-earth:lithosphere}{स्थलमंडल} से हल्का है (\cref{ch:b1:heat-budget})। महाद्वीपों के नीचे, इसके विपरीत, \omterm{def:b1:layered-earth:layers}{भूपर्पटी} की मोटाई उच्चावच के साथ बदलती देखी जाती है।
\end{solution}

\begin{solution}{exo:b1:isostasy-gravity:6}
(क) स्टेशन 1: $g - \gamma = 9.807751 - 9.808911 = \qty{-0.001160}{m.s^{-2}} =
\qty{-116.0}{\milli\gal}$; \omterm{def:b1:isostasy-gravity:anomalies}{मुक्त-वायु सुधार} $0.3086\times350 =
\qty{108.0}{\milli\gal}$, अतः $\Delta g_{\mathrm{FA}} = \qty{-8.0}{\milli\gal}$;
\omterm{def:b1:isostasy-gravity:anomalies}{बूगेर सुधार} $0.112\times350 = \qty{39.2}{\milli\gal}$, अतः
$\Delta g_{\mathrm B} = \qty{-47.2}{\milli\gal}$।
स्टेशन 2: $g - \gamma = \qty{-535.9}{\milli\gal}$; \omterm{def:b1:isostasy-gravity:anomalies}{मुक्त-वायु सुधार}
$0.3086\times1850 = \qty{570.9}{\milli\gal}$, अतः $\Delta g_{\mathrm{FA}} =
\qty{+35.0}{\milli\gal}$; \omterm{def:b1:isostasy-gravity:anomalies}{बूगेर सुधार} $0.112\times1850 =
\qty{207.2}{\milli\gal}$, अतः $\Delta g_{\mathrm B} = \qty{-172.2}{\milli\gal}$।
(ख) स्टेशन 2 पर \omterm{def:b1:isostasy-gravity:anomalies}{बूगेर विसंगति} प्रबल रूप से ऋणात्मक है: पर्वतों के नीचे द्रव्यमान की कमी है, हल्की \omterm{def:b1:layered-earth:layers}{भूपर्पटी} की एक जड़। \omterm{def:b1:isostasy-gravity:anomalies}{मुक्त-वायु विसंगति}, \omterm{def:b1:isostasy-gravity:anomalies}{बूगेर सुधार} की तुलना में छोटी, दिखाती है कि कमी पर्वतों को लगभग निरस्त करती है: वे बहुत हद तक प्रतिपूरित हैं (लगभग $172/207$, या \qty{83}{\percent}, यदि पट्टी प्रतिरूप पूरी शृंखला के लिए सही है)।
\end{solution}

\begin{solution}{exo:b1:isostasy-gravity:7}
(क) प्रस्ताव के सिद्ध की तरह, दूरी $z$ पर एक पतली परत
$G\rho\,2\pi z\,h\int_0^a s\,(s^2 + z^2)^{-3/2}\,\mathrm ds$ का योगदान देती है, अब समाकल को $s = a$ पर रोककर:
$\int_0^a s\,(s^2 + z^2)^{-3/2}\,\mathrm ds = 1/z - 1/\sqrt{z^2 + a^2}$ है। अतः
$g = 2\pi G\rho h\,\bigl(1 - z/\sqrt{z^2 + a^2}\bigr)$, जो $a \to \infty$ पर पट्टी के मान की ओर जाता है। (ख) $z/\sqrt{z^2 + a^2} = 0.1$ से
$\sqrt{z^2 + a^2} = 10z$, $a = \sqrt{99}\,z = \qty{9.95}{km}$, $z = \qty{1}{km}$ के लिए। (ग) स्टेशन के नीचे अपनी गहराई के लगभग दस गुना के भीतर की चट्टान आकर्षण का \qty{90}{\percent} देती है: भूभाग सुधार को हर स्टेशन के चारों ओर केवल उसके नीचे नहीं, बल्कि कुछ किलोमीटर से दसियों किलोमीटर त्रिज्या तक के उच्चावच की आवश्यकता है।
\end{solution}

\begin{solution}{exo:b1:isostasy-gravity:8}
(क) $r' = d(\rho_c - \rho_w)/(\rho_m - \rho_c) = 4\times1770/500 =
\qty{14.16}{km}$, और \omterm{def:b1:layered-earth:layers}{भूपर्पटी} $35 - 4 - 14.16 = \qty{16.84}{km}$ मोटी है।
(ख) एक ही भूपर्पटी-घनत्व वाला एयरी प्रतिरूप महासागरीय \omterm{def:b1:layered-earth:layers}{भूपर्पटी} को तीन गुने से अधिक मोटी बताता है। महासागरीय \omterm{def:b1:layered-earth:layers}{भूपर्पटी} पतली की गई महाद्वीपीय \omterm{def:b1:layered-earth:layers}{भूपर्पटी} नहीं है: वह एक भिन्न, सघनतर चट्टान है, और पृथ्वी के दो स्तर दो प्रकार की \omterm{def:b1:layered-earth:layers}{भूपर्पटी} हैं, जैसा \omterm{def:b1:isostasy-gravity:hypsometry}{हिप्सोमेट्रिक वक्र} सुझाता है।
\end{solution}

\begin{solution}{exo:b1:isostasy-gravity:9}
(क) एक चाप-सेकंड $\pi/(180\times3600) = 1/\num{206265}$ रेडियन है:
\qty{5.236}{''} $= \qty{2.538e-5}{rad}$ और \qty{15.885}{''} $=
\qty{7.701e-5}{rad}$। (ख) $a_h = g\theta$: $9.79\times\num{2.538e-5} =
\qty{2.485e-4}{m.s^{-2}} = \qty{24.9}{\milli\gal}$ प्रेक्षित, और
$9.79\times\num{7.701e-5} = \qty{75.4}{\milli\gal}$ दृश्य पर्वतों से परिकलित। (ग) $1 - 5.236/15.885 = 0.670$: प्रतिपूर्ति ने दृश्य पर्वतों के आकर्षण का लगभग दो-तिहाई निरस्त किया।
\end{solution}

\begin{solution}{exo:b1:isostasy-gravity:10}
(क) उसके बाहर, एक गोला अपने केंद्र पर एक बिंदु-द्रव्यमान की तरह आकर्षित करता है, केंद्र की रेखा के अनुदिश $GM/(x^2 + z^2)$ के साथ; ऊर्ध्वाधर घटक गुणा कोसाइन $z/\sqrt{x^2 + z^2}$, जो $GMz/(x^2 + z^2)^{3/2}$ देता है। (ख) $M =$ $\tfrac43\pi(\qty{2000}{m})^3\times(\qty{-200}{kg.m^{-3}}) =$
$\qty{-6.70e12}{kg}$, और केंद्र के ऊपर $GM/z^2 = \num{6.674e-11}\times
(\num{-6.70e12})/(4000)^2 = \qty{-2.80e-5}{m.s^{-2}} = \qty{-2.80}{\milli\gal}$। (ग) $(1 + x^2/z^2)^{-3/2} = 1/2$ से, अर्थात् $x = z\sqrt{2^{2/3} - 1} = 0.766\,z =
\qty{3.07}{km}$। (घ) परिच्छेदिका की अर्ध-उच्चिका-अर्ध-चौड़ाई गहराई देती है: $z = 1.305\,x_{1/2}$। गहरा पिंड चौड़ी, चपटी विसंगति देता है।
\end{solution}

\begin{solution}{exo:b1:isostasy-gravity:11}
(क) $V = \tfrac13\pi R^2 H = \tfrac13\pi(\qty{15e3}{m})^2(\qty{3e3}{m}) =
\qty{7.07e11}{m^3}$, और $\Delta m = (2800 - 1030)V = \qty{1.25e15}{kg}$।
(ख) द्रव्यमान-केंद्र तल से \qty{0.75}{km} ऊपर है, समुद्र-सतह के नीचे $d = 5 - 0.75 = \qty{4.25}{km}$ पर। वहाँ विभव-विसंगति $T = -G\Delta m/d$ है, अतः $N = G\Delta m/(\gamma d) =
\num{6.674e-11}\times\num{1.25e15}/(9.81\times4250) = \qty{2.0}{m}$।
(ग) शंकु का अधिकांश द्रव्यमान उसके चौड़े निचले भाग में है, शिखर के ऊपर के बिंदु से \qty{5}{km} से \qty{16}{km} तक: द्रव्यमान पर $1/r$ का माध्य $1/\qty{4.25}{km}$ से छोटा है, और बिंदु-द्रव्यमान विभव को अधिक आँकता है। पतली चकतियों को जोड़ने पर लगभग \qty{1.1}{m} आता है।
(घ) नीचे प्रकाश \omterm{def:b1:layered-earth:layers}{भूपर्पटी} का एक मूल उतने ही द्रव्यमान की कमी को और गहराई में रखता; दोनों विभव लगभग निरस्त हो जाते, और उभार बहुत छोटा होता, पर शून्य नहीं, क्योंकि कमी अधिकता से सतह से दूर है।
\end{solution}

\begin{solution}{exo:b1:isostasy-gravity:12}
(क) ऊपर उठने की दर $-\mathrm dw/\mathrm dt = w_0\,e^{-t/\tau}/\tau = w/\tau$ है। (ख) $w = 818\,e^{-10000/4000} = 818\times0.0821 = \qty{67.1}{m}$, और दर
$\qty{67.1}{m}/\num{4000}\,\unit{yr} = \qty{16.8}{mm.yr^{-1}}$ है, मापे गए मान से डेढ़ गुने से अधिक। (ग) $\tau = 3100$ वर्ष के साथ, $w =
818\,e^{-3.226} = \qty{32.5}{m}$ और दर \qty{10.5}{mm.yr^{-1}}; $\tau = \num{68000}$ वर्ष के साथ, $w = 818\,e^{-0.147} = \qty{706}{m}$ और दर
\qty{10.4}{mm.yr^{-1}} है। दोनों बहुत भिन्न शेष अवनमन सुझाते हैं। एक अप्रतिपूरित अवनमन प्रावार-चट्टान की द्रव्यमान-कमी है: एक विस्तृत क्षेत्र पर यह लगभग $-2\pi G\rho_m w$ की \omterm{def:b1:isostasy-gravity:anomalies}{मुक्त-वायु विसंगति} देता है, अर्थात्
$-0.0419\times3.3\times706 \approx \qty{-98}{\milli\gal}$ धीमी स्थिति के लिए, और तेज़ स्थिति के लिए \qty{-4}{\milli\gal}। एक गुरुत्व सर्वेक्षण (या पुराने तटों की ऊँचाई) उन्हें अलग करता है। (घ) समय-स्थिरांक मापता है कि अवनमन के नीचे \omterm{def:b1:layered-earth:layers}{प्रावार} कितनी तेज़ी से वापस बहता है: उसकी श्यानता।
\end{solution}

\begin{solution}{pb:b1:isostasy-gravity:1}
\textbf{1.} $r = h\rho_c/(\rho_m - \rho_c) = 4\times2700/600 = \qty{18}{km}$.

\textbf{2.} \omterm{def:b1:layered-earth:layers}{भूपर्पटी} $35 + 4 + 18 = \qty{57}{km}$ मोटी है; उसका आधार, \omterm{def:b1:layered-earth:discontinuity}{मोहो}, समुद्र-तल से नीचे $35 + 18 = \qty{53}{km}$ पर है।

\textbf{3.} पठार: \qty{57}{km} \omterm{def:b1:layered-earth:layers}{भूपर्पटी} और $60 - 53 = \qty{7}{km}$ \omterm{def:b1:layered-earth:layers}{प्रावार}, $\num{57000}\times2700 + \num{7000}\times3300 = \num{1.539e8} + \num{2.31e7} =
\qty{1.770e8}{kg.m^{-2}}$। निचला मैदान: $\num{35000}\times2700 + \num{25000}\times3300
= \num{9.45e7} + \num{8.25e7} = \qty{1.770e8}{kg.m^{-2}}$। वे बराबर हैं, जैसा \cref{prop:b1:isostasy-gravity:equal-mass} माँगता है।

\textbf{4.} $p = g\times\text{प्रति एकाई क्षेत्रफल द्रव्यमान} = 9.81\times\num{1.770e8} =
\qty{1.74e9}{Pa}$, दोनों स्तंभों के नीचे लगभग \qty{1.7}{GPa}।

\textbf{5.} मूल को अपने छोटे घनत्व से पठार के द्रव्यमान की भरपाई करनी है: मूल का प्रत्येक मीटर $\rho_m - \rho_c = \qty{600}{kg.m^{-3}}$ द्रव्यमान हटाता है, और पठार का प्रत्येक मीटर \qty{2700}{kg.m^{-3}} जोड़ता है, अतः $2700/600 = 4.5$ मीटर मूल चाहिए। \qty{2900}{kg.m^{-3}} की \omterm{def:b1:layered-earth:layers}{भूपर्पटी} के साथ अनुपात $2900/400 = 7.25$ होगा: अंतर जितना छोटा, मूल उतना गहरा।

\textbf{6.} \omterm{def:b1:layered-earth:discontinuity}{मोहो}, सीस्मिक गति में तीखी वृद्धि, पठार के किनारे के नीचे गहराता हुआ पाया जाएगा, निचले मैदान के नीचे \qty{35}{km} से पठार के नीचे लगभग \qty{53}{km} तक।

\textbf{7.} $\rho = 2700\times100/104 = \qty{2596}{kg.m^{-3}}$.

\textbf{8.} $\num{104000}\times2596.2 = \qty{2.700e8}{kg.m^{-2}}$ और
$\num{100000}\times2700 = \qty{2.700e8}{kg.m^{-2}}$।

\textbf{9.} $2700 - 2596 = \qty{104}{kg.m^{-3}}$, या $104/2700 =
\qty{3.85}{\percent}$।

\textbf{10.} $\Delta\rho/\rho = \alpha\,\Delta T$ से $\Delta T =
0.0385/\num{3e-5} = \qty{1280}{K}$ पूरे \qty{104}{km} स्तंभ पर औसतन। वह स्तंभ के बड़े भाग को पिघला देगा: अकेली ऊष्मा इस पठार को इतना हल्का नहीं कर सकती, और प्रैट की व्यवस्था यहाँ मुख्य नहीं हो सकती।

\textbf{11.} प्रैट के स्तंभों में कोई मूल नहीं है: \omterm{def:b1:layered-earth:discontinuity}{मोहो} पठार के नीचे लगभग \qty{35}{km} पर रहता है, निचले मैदान के नीचे की तरह।

\textbf{12.} \qty{52}{km} एयरी के \qty{53}{km} के निकट है और प्रैट के \qty{35}{km} से दूर: पठार मुख्य रूप से \omterm{def:b1:layered-earth:layers}{भूपर्पटी} के एक मूल से प्रतिपूरित है।

\textbf{13.} $0.0419\times2.7\times4000 = \qty{452.5}{\milli\gal}$.

\textbf{14.} अप्रतिपूरित (\cref{prop:b1:isostasy-gravity:broad}):
$\Delta g_{\mathrm{FA}} \approx \qty{+452.5}{\milli\gal}$,
$\Delta g_{\mathrm B} \approx 0$।

\textbf{15.} प्रतिपूरित: $\Delta g_{\mathrm{FA}} \approx 0$ और
$\Delta g_{\mathrm B} \approx \qty{-452.5}{\milli\gal}$। मूल के लुप्त द्रव्यमान से
$-0.0419\times0.6\times\num{18000} = \qty{-452.5}{\milli\gal}$ मिलता है, वही, क्योंकि $\Delta\rho\,r = \rho_c h$।

\textbf{16.} $\sin^2\ang{30} = 0.25$ और $\sin^2\ang{60} = 0.75$:
$\gamma = 9.780327\times(1 + 0.0013256 - 0.0000044) = \qty{9.793249}{m.s^{-2}}$।
तब $g - \gamma = 9.781085 - 9.793249 = \qty{-1216.4}{\milli\gal}$; \omterm{def:b1:isostasy-gravity:anomalies}{मुक्त-वायु सुधार}
$0.3086\times4000 = \qty{1234.4}{\milli\gal}$, अतः
$\Delta g_{\mathrm{FA}} = \qty{+18.0}{\milli\gal}$, और
$\Delta g_{\mathrm B} = 18.0 - 452.5 = \qty{-434.5}{\milli\gal}$।

\textbf{17.} $434.5/452.5 = 0.96$: पठार लगभग \qty{96}{\percent} प्रतिपूरित है (छोटी धनात्मक \omterm{def:b1:isostasy-gravity:anomalies}{मुक्त-वायु विसंगति} अप्रतिपूरित \qty{4}{\percent} है)।

\textbf{18.} नहीं। एक विस्तृत पठार के केंद्र के निकट हर असामान्य परत एक पट्टी की तरह कार्य करती है, जिसका आकर्षण उसकी गहराई पर निर्भर नहीं करता: केवल प्रति एकाई क्षेत्रफल द्रव्यमान-कमी मायने रखती है, और वह दोनों प्रतिरूपों में समान है, $\rho_c h$। वे केवल किनारों के निकट भिन्न हैं, जहाँ प्रतिपूर्ति करने वाले द्रव्यमानों की गहराई विसंगति का आकार बदलती है।

\textbf{19.} किसी भी ओर किनारे से दूर \omterm{def:b1:isostasy-gravity:anomalies}{मुक्त-वायु विसंगति} लगभग शून्य है। किनारे के भीतर कुछ किलोमीटर, स्टेशन के नीचे का पठार निकट है और केवल एक ओर दूर तक फैला है: वह अब भी लगभग अपने पूरे पट्टी-मान से खींचता है। मूल, \qty{35}{km} या उससे अधिक नीचे, वहाँ से लगभग सीधे नीचे समाप्त होता एक पिंड दिखता है: वह अपने प्रभाव के आधे से थोड़ा ही अधिक खींचता है। अंतर एक धनात्मक शिखर है। कुछ किलोमीटर बाहर, पठार बहुत कम खींचता है, जबकि मूल अब भी अपने प्रभाव के लगभग आधे से खींचता है: एक ऋणात्मक गर्त। परिच्छेदिका हर किनारे के पार एक धनात्मक-ऋणात्मक जोड़ी है, जैसा अध्याय के चित्र में।

\textbf{20.} \omterm{def:b1:layered-earth:layers}{भूपर्पटी} अब $57 - 1 = \qty{56}{km}$ मोटी है, निचले मैदान से $\qty{21}{km}$ अधिक। साम्य में यह अधिकता $h' + r'$ में बँटती है, $r' = 4.5\,h'$ के साथ: $h' = 21/5.5 = \qty{3.818}{km}$ और $r' = \qty{17.18}{km}$।

\textbf{21.} पठार अब \qty{3818}{m} पर खड़ा है।

\textbf{22.} सतह पर अब जो चट्टान है, वह पुरानी सतह से \qty{1}{km} नीचे थी, \qty{3000}{m} पर; अब वह \qty{3818}{m} पर है: वह \qty{818}{m} ऊपर उठ गई है, अर्थात् $1000\times\rho_c/\rho_m$।

\textbf{23.} हर क्षरित किलोमीटर सतह को $1 - 2700/3300 =
0.182$ किलोमीटर नीचे लाता है, अतः उसे \qty{4}{km} नीचे लाने में
$4/0.182 = \qty{22}{km}$ क्षरण लगता है, यानी पूरी अतिरिक्त मोटाई $h + r$। एक पुराने, घिसे पर्वत-पट्टे में सतह की चट्टानें कभी लगभग बीस किलोमीटर गहरी थीं: पट्टे की गहरी जड़ें उजागर हो गई हैं।

\textbf{24.} सतह $1000 - 818$ से नीचे आती है: क्षरित चट्टान के हर किलोमीटर के लिए सतह की शुद्ध गिरावट
\textbf{\qty[text-series-to-math=true]{182}{m}}।
\end{solution}
